\Article{20}
It remains to determine the nature of the equilibrium shape from the theorem in Article 18. The central task is to work out in general the variation of the expression $W$ when the shape of the region filled with fluid undergoes an arbitrary infinitely small change. But since the calculus of variations of double integrals is as yet little developed for the case in which the limits must also be treated as variables, we must go a little deeper into this delicate question.

We shall take the surface that separates the region $s$ from the rest of space and consider its part $U$, supposing each point of that part to be determined by three coordinates $x$, $y$, $z$, the third of which is the distance from an arbitrary horizontal plane. Thus $z$ can be regarded as a function of the indeterminates $x$, $y$. We shall denote its partial differentials in the usual way, with the grouping marks omitted, by
\[
\frac{\mathrm{d}z}{\mathrm{d}x}.\mathrm{d}x,\qquad\frac{\mathrm{d}z}{\mathrm{d}y}.\mathrm{d}y
\]
At any point of the surface we imagine a straight line normal to the surface and directed outward with respect to the region $s$. We denote the cosines of the angles between this normal and the straight lines parallel to the coordinate axes $x$, $y$, $z$ by $\xi$, $\eta$, $\zeta$. Thus
\[
\begin{gathered}
\xi\xi+\eta\eta+\zeta\zeta=1\\
\frac{\mathrm{d}z}{\mathrm{d}x}=-\frac{\xi}{\zeta},\qquad\frac{\mathrm{d}z}{\mathrm{d}y}=-\frac{\eta}{\zeta}
\end{gathered}
\]
The boundary of the surface $U$ will be a line that closes on itself, which we denote by $P$, and, supposing it to be \SourceLeaf{62}traced by a continuous motion, we shall always take its elements $\mathrm{d}P$ to be positive, just like the surface elements $\mathrm{d}U$. We denote the cosines of the angles that the direction of the element $\mathrm{d}P$ makes with the coordinate axes $x,y,z$ by $X,Y,Z$. To leave no ambiguity about the sense of that direction, we choose it so that it comes first, the direction normal to the element $\mathrm{d}P$ that is tangent to the surface $U$ and directed inward with respect to it comes second, and the normal to the surface directed outward with respect to the region $s$ comes third. In this order the three straight lines have the same relative orientation as the coordinate axes $x,y,z$. It is readily seen (cf. \textit{Disquiss. gen. circa superficies curvas} art. 2) that the cosines of the angles between that second direction and the coordinate axes $x,y,z$ are respectively
\[
\eta^0Z-\zeta^0Y,\qquad \zeta^0X-\xi^0Z,\qquad \xi^0Y-\eta^0X
\]
if $\xi^0,\eta^0,\zeta^0$ are the values of $\xi,\eta,\zeta$ at the point of the element $\mathrm{d}P$.

\Article{21}
With these preliminaries in place, let us suppose that the surface $U$ undergoes an arbitrary infinitely small change. If it were enough to consider only changes for which the boundary $P$ always remained unchanged or at least remained on the same vertical surface, we should evidently have to introduce a variation only in the third coordinate $z$, which would make the problem much easier. But since we have to treat the problem with the greatest generality, in such an approach taking account of the variability of the limits would lead to awkward complications that would spoil the elegance of the treatment. It is therefore preferable to subject all three coordinates to variation from the outset. We shall therefore picture the matter in this way. For each point of the surface whose coordinates are $x,y,z$ we substitute another whose coordinates are $x+\delta x$, $y+\delta y$, $z+\delta z$. Here $\delta x,\delta y,\delta z$ may be regarded as undetermined functions of $x,y$, but their values remain infinitely small. Let us now examine the variations of the individual elements of the expression $W$, beginning with the variation of the element $\mathrm{d}U$.

Imagine a triangular element of the surface $U$, denoted by $\mathrm{d}U$, between the points with coordinates
\[
\begin{array}{lll}
x,&y,&z\\
x+\mathrm{d}x,&y+\mathrm{d}y.&z+\frac{\mathrm{d}z}{\mathrm{d}x}.\mathrm{d}x+\frac{\mathrm{d}z}{\mathrm{d}y}.\mathrm{d}y\\
x+\mathrm{d}'x,&y+\mathrm{d}'y,&z+\frac{\mathrm{d}z}{\mathrm{d}x}.\mathrm{d}'x+\frac{\mathrm{d}z}{\mathrm{d}y}.\mathrm{d}'y
\end{array}
\]
\SourceLeaf{63}Twice the area of this triangle is found by familiar principles to be
\[
=(\mathrm{d}x.\mathrm{d}'y-\mathrm{d}y.\mathrm{d}'x)\BareRoot{[1+(\frac{\mathrm{d}z}{\mathrm{d}x})^2+(\frac{\mathrm{d}z}{\mathrm{d}y})^2]}
\]
if we suppose, as we may, that $\mathrm{d}x.\mathrm{d}'y-\mathrm{d}y.\mathrm{d}'x$ is a positive quantity.

On the varied surface we shall have three other points in place of those points. Their coordinates are as follows.
\[
\begin{array}{ll}
\text{of the first point}&x+\delta x,\quad y+\delta y,\quad z+\delta z
\end{array}
\]
of the second point
\[
\begin{gathered}
x+\mathrm{d}x+\delta x+\frac{\mathrm{d}\delta x}{\mathrm{d}x}.\mathrm{d}x+\frac{\mathrm{d}\delta x}{\mathrm{d}y}.\mathrm{d}y\\
y+\mathrm{d}y+\delta y+\frac{\mathrm{d}\delta y}{\mathrm{d}x}.\mathrm{d}x+\frac{\mathrm{d}\delta y}{\mathrm{d}y}.\mathrm{d}y\\
z+\frac{\mathrm{d}z}{\mathrm{d}x}.\mathrm{d}x+\frac{\mathrm{d}z}{\mathrm{d}y}.\mathrm{d}y+\delta z+\frac{\mathrm{d}\delta z}{\mathrm{d}x}.\mathrm{d}x+\frac{\mathrm{d}\delta z}{\mathrm{d}y}.\mathrm{d}y
\end{gathered}
\]
of the third point
\[
\begin{gathered}
x+\mathrm{d}'x+\delta x+\frac{\mathrm{d}\delta x}{\mathrm{d}x}.\mathrm{d}'x+\frac{\mathrm{d}\delta x}{\mathrm{d}y}.\mathrm{d}'y\\
y+\mathrm{d}'y+\delta y+\frac{\mathrm{d}\delta y}{\mathrm{d}x}.\mathrm{d}'x+\frac{\mathrm{d}\delta y}{\mathrm{d}y}.\mathrm{d}'y\\
z+\frac{\mathrm{d}z}{\mathrm{d}x}.\mathrm{d}'x+\frac{\mathrm{d}z}{\mathrm{d}y}.\mathrm{d}'y+\delta z+\frac{\mathrm{d}\delta z}{\mathrm{d}x}.\mathrm{d}'x+\frac{\mathrm{d}\delta z}{\mathrm{d}y}.\mathrm{d}'y
\end{gathered}
\]
Twice the area of the triangle between these points is found by the same method to be
\[
=(\mathrm{d}x.\mathrm{d}'y-\mathrm{d}y.\mathrm{d}'x)\BareRoot{N}
\]
if for brevity $N$ denotes the sum
\[
\begin{aligned}
&[(1+\frac{\mathrm{d}\delta x}{\mathrm{d}x})(1+\frac{\mathrm{d}\delta y}{\mathrm{d}y})-\frac{\mathrm{d}\delta x}{\mathrm{d}y}.\frac{\mathrm{d}\delta y}{\mathrm{d}x}]^2\\
+&[(1+\frac{\mathrm{d}\delta x}{\mathrm{d}x})(\frac{\mathrm{d}z}{\mathrm{d}y}+\frac{\mathrm{d}\delta z}{\mathrm{d}y})-\frac{\mathrm{d}\delta x}{\mathrm{d}y}(\frac{\mathrm{d}z}{\mathrm{d}x}+\frac{\mathrm{d}\delta z}{\mathrm{d}x})]^2\\
+&[(1+\frac{\mathrm{d}\delta y}{\mathrm{d}y})(\frac{\mathrm{d}z}{\mathrm{d}x}+\frac{\mathrm{d}\delta z}{\mathrm{d}x})-\frac{\mathrm{d}\delta y}{\mathrm{d}x}(\frac{\mathrm{d}x}{\mathrm{d}y}+\frac{\mathrm{d}\delta z}{\mathrm{d}y})]^2
\end{aligned}
\]
Expanding and discarding quantities of the second order, we find
\[
\BareRoot{N}=\BareRoot{[1+(\frac{\mathrm{d}z}{\mathrm{d}x})^2+(\frac{\mathrm{d}z}{\mathrm{d}y})^2]}.\left[1+\frac{L}{1+(\frac{\mathrm{d}z}{\mathrm{d}x})^2+(\frac{\mathrm{d}z}{\mathrm{d}y})^2}\right]
\]
if for brevity $L$ denotes the sum
\SourceLeaf{64}\[
\begin{aligned}
&\frac{\mathrm{d}\delta x}{\mathrm{d}x}[1+(\frac{\mathrm{d}z}{\mathrm{d}y})^2]-\frac{\mathrm{d}\delta x}{\mathrm{d}y}.\frac{\mathrm{d}z}{\mathrm{d}x}.\frac{\mathrm{d}z}{\mathrm{d}y}-\frac{\mathrm{d}\delta y}{\mathrm{d}x}.\frac{\mathrm{d}z}{\mathrm{d}x}.\frac{\mathrm{d}z}{\mathrm{d}y}\\
&+\frac{\mathrm{d}\delta y}{\mathrm{d}y}[1+(\frac{\mathrm{d}z}{\mathrm{d}x})^2]+\frac{\mathrm{d}\delta z}{\mathrm{d}x}.\frac{\mathrm{d}z}{\mathrm{d}x}+\frac{\mathrm{d}\delta z}{\mathrm{d}y}.\frac{\mathrm{d}z}{\mathrm{d}y}
\end{aligned}
\]
The ratio of the first triangle to the second is therefore as $1$ to
\[
1+\frac{L}{1+(\frac{\mathrm{d}z}{\mathrm{d}x})^2+(\frac{\mathrm{d}z}{\mathrm{d}y})^2}
\]
and is thus independent of the shape of the triangle $\mathrm{d}U$. The result is
\[
\delta\mathrm{d}U=\frac{L\mathrm{d}U}{1+(\frac{\mathrm{d}z}{\mathrm{d}x})^2+(\frac{\mathrm{d}z}{\mathrm{d}y})^2}
\]
or, in explicit terms,
\[
\begin{aligned}
\delta\mathrm{d}U=\mathrm{d}U\big(&\frac{\mathrm{d}\delta x}{\mathrm{d}x}(\eta\eta+\zeta\zeta)-\frac{\mathrm{d}\delta x}{\mathrm{d}y}.\xi\eta-\frac{\mathrm{d}\delta y}{\mathrm{d}x}.\xi\eta\\
&+\frac{\mathrm{d}\delta y}{\mathrm{d}y}(\xi\xi+\zeta\zeta)-\frac{\mathrm{d}\delta z}{\mathrm{d}x}.\xi\zeta-\frac{\mathrm{d}\delta z}{\mathrm{d}y}.\eta\zeta\big)
\end{aligned}
\]

